\documentclass{jsarticle}
\usepackage{amsmath}
\usepackage{amssymb}
\begin{document}
\title{数学と数がオイラーの定数から生まれた　}
\author{Masaaki Yamaguchi}
\date{}
\maketitle

   $$ 2^2 = e^{x \log x} = 4 $$
   $$ 3^3 = e^{x \log x} = 27 $$
   $$ 4^4 = e^{x \log x} = 256 $$
   $$ 5^5 = e^{x \log x} = 3125 $$


    $$ y = x \log x, \log x \to (\log x)^{-1}  $$
    $$ {1 \over 2} + iy = {{x \log x} \over \log x}$$


   $$ x^{{1 \over 2} + iy} = e^{x \log x} $$
   $$ {1 \over 2} + iy = \log_{x}e^{x \log x} = \log_{x} 4, \\
   \log_{x} 27, \log_{x} 256, \log_{x} 3125 $$
   $$ y = e^{x \log x} = \sqrt{a} $$
   $$ e^{e^{x \log x}} = a, e^{(x\log x)^2} $$
   $$ x^2 = \pm a, \lim_{n \to \infty} (x - y) = e^{x \log x} $$
   $$ \int {1 \over (x \log x)}dx = i \int x \log x dx 
   - \int {1 \over (x \log x)}dx $$
   $$ y = x \log x, \log x \to (\log x)^{-1}  $$
   $$ ||ds^2|| ~ 8 \pi G\left({p \over c^3} + {V \over S}\right)$$
   と宇宙の中の１種の原子をみつける正確さがこの式と、


   $$ y = {{x \log x} \over (\log x)} = x $$
   と、$x \log x = a$から ${a \over (\log x)} \to x $と x を抜き取る。
   このx　を見つけるのに　$x^{{1 \over 2} + iy} = e^{x \log x} $
   ${1 \over 2} + iy= {{x \log x} \over (\log x)} = x$として
   このｘを見つける式がゼータ関数である。


   ゼータ関数は、量子暗号にもなっていることと、この式自体が公開鍵暗号文
    にもなっている。

   $$ {1 \over 2} + iy = {{x \log x} \over \log x}$$
   この式が一次独立であるためには。
   　$$x = {1 \over 2}, iy = 0　$$
   がゼータ関数となる必要十分条件でもある。

   $$ \int C dx_m = 0 $$
   $$ {d \over df} \int C dx_m = 0^{0^{'}} $$
   $$ = e^{x \log x} $$
   と標数０の体の上の代数多様体でもあり、このオイラーの定数からの
   大域的微分多様体から数が生まれた。

   $$ H \Psi = \bigoplus {{({i \hbar}^{\nabla})}^{\oplus L}} $$
   $$ \bigoplus a^f x^{1 - f}[I_m] $$
   $$= \int e^x x^{1 - t}dx_m $$
   $$ = e^{x \log x} $$



    アメリカ大統領を統計で選ぶ選挙は、reco level理論がゼータ関数
    として機能する遷移エネルギーの安定軌道をある集団ｘに対数$\log x$
    の組み合わせとして、指数の巨大確率を対数の個数とする
    この大統領の素質としての$x^n $集団の共通の思考がｎとなるこのｎがどのくらいの
    エントロピー量かを$H = - K p \log p $ が
    表している。


    $$ \int \Gamma(\gamma)^{'}dx_m = (e^f + e^{-f}) \ge (e^{f} - e^{-f}) $$
    $$(e^f + e^{-f}) \ge (e^{f} - e^{-f}) $$ 
    この方程式はブラックホールのシュバルツシルト半径から
    $$ (e^f + e^{-f})(e^f - e^{-f}) = 0 $$
    $$ {d \over df}F \cdot \int C dx_m \ge 0 $$
    $$ y = f(g(x)^{'})dx = \int f(x)^{'}g(x)^{'}dx $$
    $$ y = f(\log x)^{'}dx = f^{'}(x){1 \over x} $$
    $$ y = {{f^{'}(x)} \over x} $$
    $$ = 2(\cos (i x \log x) - i \sin (i x \log x)) $$
    $$ = C $$
    $$ c = f(x)\cdot \log x, dx_m = (\log x)^{-1}  $$
    $$ C = {d \over \gamma}\Gamma = \bigoplus (i \hbar^{\nabla})^{\oplus L} $$
    となり、ヴェイユ予想の式からも導かれる。
    $$ = 2(\cos (i x \log x) - i \sin (i x \log x)) $$
    $$ ~ e^{\theta} $$
    $$ {d \over \gamma}\Gamma = \Gamma^{{\gamma}^{'}} $$
    $$ = \int \Gamma(\gamma)^{'}dx_m $$
    $$ y = f(x)\log x, y^{'} = f^{'}(x) + {f^{'}(x) \over x} $$

    $$ \sin (\log x)^{'} dx = \cos (\log x)\cdot {1 \over x} $$
    $$ \sin{y \over x} = \sin \vec{u} = a + t \sin \vec{u} $$
    $$ ~ i, -i, 2i, -2i $$
    $$ \lim_{n \to \infty}(f(b) - f(a)) = f^{'}(c)(b - a) $$
    $$ \sin (\log x)^{'}dx = \cos (\log x)\cdot {1 \over x} $$
    $$ y = f(x)\log x, y^{'} = f^{'}(x) + {f^{'}(x) \over x} $$
    $$ \sin (\log x)^{'} dx = \cos (\log x)\cdot {1 \over x} $$
    $$ \sin{y \over x} =  2(\cos (i x \log x) - i \sin (i x \log x)) $$
    $$ = e^{\theta} $$

\end{document}
